% !TeX program = xelatex
\documentclass[UTF8,a4paper,12pt]{ctexart}
\usepackage{geometry}
\usepackage{listings}
\usepackage{xcolor}

\geometry{a4paper, top=0.4in, bottom=1in, left=1in, right=1in}

\lstset{
    language=C++,
    basicstyle=\ttfamily\small,
    keywordstyle=\color{blue},
    commentstyle=\color{gray},
    stringstyle=\color{orange},
    numbers=left,
    numberstyle=\tiny\color{gray},
    frame=single,
    breaklines=true
}

\title{四则运算}
\author{GESP 2级}
\date{}

\begin{document}

\maketitle

\section*{四则运算 (Arithmetic)}

\subsection*{整数运算 (Integer Operations)}

\begin{lstlisting}
int a = 10, b = 3;
cout << a + b << endl;  // 13（加法）
cout << a - b << endl;  // 7（减法）
cout << a * b << endl;  // 30（乘法）
cout << a / b << endl;  // 3（整数除法）
cout << a % b << endl;  // 1（取余）
\end{lstlisting}

\subsection*{浮点数运算 (Floating-point Operations)}

\begin{lstlisting}
double x = 10.5, y = 2.5;
cout << x + y << endl;  // 13.0（加法）
cout << x - y << endl;  // 8.0（减法）
cout << x * y << endl;  // 26.25（乘法）
cout << x / y << endl;  // 4.2（浮点除法）
\end{lstlisting}

\section*{词汇表 (Glossary)}

\begin{center}
\begin{tabular}{|l|l|p{10cm}|}
\hline
\textbf{英文名词} & \textbf{中文名词} & \textbf{解释} \\
\hline
Integer & 整数 & 不包含小数部分的数字，如 10、-5、0。在C++中用 \texttt{int} 类型表示。 \\
\hline
Floating-point & 浮点数 & 包含小数部分的数字，如 10.5、3.14。在C++中用 \texttt{double} 或 \texttt{float} 表示。 \\
\hline
Addition & 加法 (+) & 两个数相加，求和。 \\
\hline
Subtraction & 减法 (-) & 两个数相减，求差。 \\
\hline
Multiplication & 乘法 (*) & 两个数相乘，求积。 \\
\hline
Division & 除法 (/) & 两个数相除，求商。整数除法会截断小数部分。 \\
\hline
Modulo & 取余 (\%) & 求两个整数相除后的余数，仅适用于整数类型。 \\
\hline
Divisibility & 整除 & 当 \texttt{a \% b == 0} 时，称 a 能被 b 整除。 \\
\hline
\end{tabular}
\end{center}

\end{document}
